\section{Explicit regional-energy parameters}

The following numerical choices control the energy of an untrimmed
candidate obtained from the corrected hard polynomial family. The soft
regional operator used to measure this energy is chosen separately.

\begin{theorem}[A fixed soft residual]
    \label{thm:soft_regional_fixed_residual}
    \lean{Real.soft_regional_fixed_residual_lt}
    \leanok
    Set $t=768$, $m=12$ and $T=384$. Then the residual
    $\rho=32e^{-(T-48)/8}+(t/m)(T/t)^m$ satisfies $0\le\rho<1/16$.
\end{theorem}
\begin{proof}\leanok
    The expression equals $32e^{-42}+1/64$. Since $e^{14}\ge15$,
    $e^{42}\ge3375$, and hence
    $\rho\le32/3375+1/64<1/16$.
\end{proof}

\begin{theorem}[The count dominates the regional norm]
    \label{thm:soft_regional_fixed_norm}
    \lean{Real.soft_regional_fixed_norm_bound}
    \leanok
    Let $\ell,D\in\N$, with $\ell\ge21$ and $D\ge50\ell^2$.
    Set $A=768(1+\log12)+4\ell+14$. Then $0\le A$ and $A+1\le D$.
\end{theorem}
\begin{proof}\leanok
    The inequality $\log12\le11$ gives
    $A+1\le9231+4\ell\le50\ell^2\le D$.
    Nonnegativity follows from $\log12\ge0$.
\end{proof}

\begin{theorem}[Precision at the same family count]
    \label{thm:soft_regional_filter_precision}
    \lean{Real.soft_regional_filter_precision_bounds}
    \leanok
    Let $D,s\in\N$, with $D\ge2$ and $s\ge1$, and set
    \begin{align}
        a      & =\frac3{32D^{12}s^5},  &
        \tau   & =\frac{a^2}{2^{20}},   &
        \Delta & =\frac1{200000D^{12}}.
    \end{align}
    Both $\tau$ and $\Delta/(1-1/4)^2$ are positive, and
    \begin{align}
        \tau                     & \le\frac1{1024(D+1)}, &
        \frac{\Delta}{(1-1/4)^2} & \le\frac1{100(D+1)}.
    \end{align}
\end{theorem}
\begin{proof}\leanok
    Since $D^{12}s^5\ge D$, one has $a\le3/(32D)$ and
    $\tau\le9/(2^{30}D^2)$. Also $D^{12}\ge D$.
    Multiply the desired inequalities by their positive denominators
    and use $D\ge2$.
\end{proof}

If the actual family count is $D=(b\ell^2)^{22\ell}$ with $b\ge50$
and $\ell\ge21$, then $D\ge50\ell^2$. These estimates therefore verify
both precision hypotheses of Lemma~\ref{lem:regional_energy_calibration}
for the same count, using one application of the improving family to
an input of viability error at most $1/4$.

The soft definition allows $t\ge m\ge2$, as satisfied here;
see the definition of soft truncation in \cite[lines 1178--1185]
{Arad2016RigorousRG}. The literal soft-generation procedure instead
forms its polynomial from the soft Hamiltonian and uses $m=4$;
see \cite[Figure \texttt{fig:generate} and lines 1377--1404]
{Arad2016RigorousRG}. The estimates above do not identify the corrected
hard polynomial family with that procedure or assert its running time.

\begin{theorem}[Calibrated precision for a terminal image]
    \label{thm:soft_terminal_energy_precision}
    \lean{Real.soft_terminal_energy_precision_bounds}
    \leanok
    Let $D\ge2$, $M\ge34$, $0\le A\le M$, and
    $0\le\Delta\le(D^{12}M^{10})^{-1}$. Set
    $\tau=(D^{24}M^{10})^{-1}$. Then
    \begin{align}
        \tau M\le\frac1{64},\qquad
        \frac{3\tau}{200}\le\frac1{1024(A+1)},\qquad
        \frac{\Delta}{(91/100)^2}\le\frac1{100(A+1)}.
    \end{align}
    These scalar bounds support the regional-energy estimate in
    \cite[line 1807]{Arad2016RigorousRG}.
\end{theorem}
\begin{proof}\leanok
    Since $M^{10}\ge M$, the two denominators are at least $64M$
    and $256M$, respectively. Also $A+1\le2M$. Thus
    $\tau M\le1/64$ and $\Delta M\le1/256$, which give all three
    estimates after multiplication by the positive denominators.
\end{proof}

\begin{theorem}[The reverse tail bounds the regional residual]
    \label{thm:soft_terminal_regional_residual}
    \lean{Real.soft_terminal_regional_residual_le}
    \leanok
    Let $g,M,\tau,t,T$ be real and $m$ a nonnegative integer.
    Suppose $0\le g\le M$, $\tau\ge0$, $\tau M\le1$, and
    \begin{align}
        3\left(32e^{-(T-156-2g)/8}
        +\frac tm\left(\frac Tt\right)^m\right)
        \le\frac{\tau g}{10^6}.
    \end{align}
    Then $32e^{-(T-48)/8}+(t/m)(T/t)^m\le1/16$.
    Division is understood with its usual total-field extension at zero.
\end{theorem}
\begin{proof}\leanok
    The exponential with shift $48$ is at most the one with shift
    $156+2g$, and $\tau g\le\tau M\le1$. The stated bound follows.
\end{proof}

\begin{theorem}[Comparing the two soft norm scales]
    \label{thm:soft_terminal_regional_norm}
    \lean{Real.soft_terminal_regional_norm_bounds}
    \leanok
    Let $t,\ell\ge0$ and $m\ge1$. Put
    $A=t(1+\log m)+4\ell+14$ and
    $M=3t(1+\log m)+4\ell+34$. Then
    $0\le A\le M$, $3t\le M$ and $34\le M$.
\end{theorem}
\begin{proof}\leanok
    Use $\log m\ge0$ and $t\log m\ge0$ in the two defining formulas.
\end{proof}
